\documentclass{beamer}
% reduce font size


\title{Spectrometry Lab: Determining Thermodynamic Properties}
\author{Patryk Kozlowski \\ \href{mailto:patryk_kozlowski@g.harvard.edu}{patryk\_kozlowski@g.harvard.edu}}
\date{Lab 4 - General Chemistry}


\begin{document}

\frame{\titlepage}

% Single Slide Content
\begin{frame}{Lab 4: Spectrophotometry and Thermodynamics}
    \tiny
    \textbf{Objective:} Use spectrophotometry to measure equilibrium concentrations at different temperatures and determine thermodynamic properties like $\Delta G^\circ$, $\Delta H^\circ$, and $\Delta S^\circ$. \\
    
    \pause
    \textbf{Reaction:} Complexation of Aluminum Ions with Xylenol Orange:
    \[
    \text{H}_4\text{Q} + \text{Al}^{3+} \rightleftharpoons \text{AlQ}^{-} + 4\text{H}^{+}
    \]
    \begin{itemize}
        \item Yellow reactant ($\text{H}_4\text{Q}$) forms red product ($\text{AlQ}^{-}$). \pause
        \item Measure absorbance at 550 nm (product absorbs more strongly at this wavelength). \pause
        \item The reaction is endothermic (heat is required to drive the reaction forward). \pause
    \end{itemize}
    
    \textbf{Key Calculations:} \pause
    \begin{itemize}
        \item Beer’s Law: \(A = \epsilon_1[\text{H}_4\text{Q}]L + \epsilon_2[\text{AlQ}^{-}]L\) \pause
        \item Find equilibrium constant \(K\) at each temperature. \pause
        \item Calculate $\Delta G^\circ$ from $K$: \(\Delta G^\circ = -RT \ln K\) \pause
        \item Use $\Delta G^\circ = \Delta H^\circ - T \Delta S^\circ$ to determine $\Delta H^\circ$ and $\Delta S^\circ$ from a plot of $\Delta G^\circ$ vs. $T$. \pause
    \end{itemize}

    \textbf{Experiment Steps:}
    \begin{itemize}
        \item Measure initial absorbance $A_i$ at 550 nm. \pause
        \item Heat the sample and record absorbance at different temperatures (90°C to 50°C). \pause
        \item Use absorbance data to find $[ \text{AlQ}^{-} ]$ and calculate $K$ at each temperature. \pause
    \end{itemize}

\end{frame}

\end{document}
